\chapter{روش پیشنهادی}
\label{ch:method}

\section{معماری سه‌مرحله‌ای}
\label{sec:architecture}

سامانه از سه مرحله متوالی ساخته شده است که جدول \ref{tab:pipeline} آن‌ها را خلاصه می‌کند. مرز میان مرحله‌ها عمدی است: هر مرحله ورودی‌اش را از فایل می‌خواند و خروجی‌اش را در فایل می‌نویسد، پس می‌شود هر کدام را جدا اجرا کرد، جدا سنجید و در صورت شکست از همان‌جا ادامه داد.

\begin{table}[ht]
\centering
\caption{سه مرحله خط لوله و کاری که هر مرحله انجام می‌دهد.}
\label{tab:pipeline}
\small
\begin{tabular}{|p{2.9cm}|p{3.5cm}|p{7.3cm}|}
\hline
\rg مرحله & \rg ورودی و خروجی & \rg کار \\
\hline
\rg ورود و پاک‌سازی & \rg سند خام به قطعه متنی & \rg استخراج متن، نرمال‌سازی، قطعه‌بندی، ثبت \lr{chunk\_sha} \\
\hline
\rg استخراج دانش & \rg قطعه به سه‌تایی & \rg فراخوانی مدل زبانی، اعتبارسنجی، رد موارد خارج از هستان‌شناسی، الزام شاهد، حذف تکراری \\
\hline
\rg بستار و بارگذاری & \rg سه‌تایی به گراف & \rg اجرای سه قاعده استنتاج، سپس \lr{MERGE} گره و یال در \lr{Neo4j} \\
\hline
\end{tabular}
\end{table}

\section{هستان‌شناسی}
\label{sec:ontology}

هستان‌شناسی داخل کد نیست و از یک فایل \lr{JSON} خوانده می‌شود. وقتی این تصمیم را گرفتیم بیشتر به‌خاطر مرتب بودن کد بود، ولی بعداً معلوم شد همین باعث می‌شود بدون یک خط بازنویسی، همان کد روی دو داده کاملاً متفاوت کار کند. جدول \ref{tab:ontology} این دو را کنار هم می‌گذارد.

\begin{table}[ht]
\centering
\caption{یک کد، دو هستان‌شناسی. برای \lr{Re-DocRED} حتی یک قاعده هم دستی نوشته نشد.}
\label{tab:ontology}
\small
\begin{tabular}{|r|c|c|c|}
\hline
داده & منبع هستان‌شناسی & نوع موجودیت & ترکیب مجاز \\
\hline
شهرداری (فارسی) & دست‌نویس & $6$ & $5$ \\
\hline
\lr{Re-DocRED} (انگلیسی) & خودکار از بخش \lr{Train} & $6$ & $654$ \\
\hline
\end{tabular}
\end{table}

آن $654$ ترکیب مجاز را یک اسکریپت از روی داده آموزشی درآورد. بارگذار \lr{Neo4j} هم همین فایل را می‌خواند، پس فیلتر نوع در هر دو سر خط لوله یکسان است.

\section{لایه بستار قاعده‌محور}
\label{sec:closure}

سه قاعده روی گراف اجرا می‌شوند:

\begin{itemize}
\item قاعده وارون: اگر $A$ در $B$ واقع باشد، آنگاه $B$ شامل $A$ است.
\item قاعده تراگذری: اگر $A$ در $B$ و $B$ در $C$ باشد، آنگاه $A$ در $C$ است.
\item قاعده زنجیره: اگر $A$ در $B$ باشد و کشور $B$ برابر $C$ باشد، آنگاه کشور $A$ هم $C$ است.
\end{itemize}

هر حقیقتی که این لایه می‌سازد، شاهد مقدمه‌های خودش را به ارث می‌برد، وگرنه قانون «هر یال باید شاهد داشته باشد» همین‌جا می‌شکست. یک تست خودکار دقیقاً همین را بررسی می‌کند.

\section{مهندسی مقیاس}
\label{sec:scale}

اجرای $500$ سند به‌صورت تک‌نخی حدود هفت ساعت طول کشیده بود. با همان سرعت، پیکره کامل بیش از یک ماه وقت می‌خواست که عملاً یعنی نشدنی. چهار تغییر این وضع را عوض کرد و هر چهارتا را روی داده واقعی اندازه گرفتیم. جدول \ref{tab:scale} نتیجه را نشان می‌دهد.

\begin{table}[ht]
\centering
\caption{چهار تغییری که اجرای پیکره کامل را ممکن کرد.}
\label{tab:scale}
\small
\begin{tabular}{|p{2.7cm}|p{4.5cm}|p{6.5cm}|}
\hline
\rg جزء & \rg تغییر & \rg اندازه‌گیری \\
\hline
\rg استخراج & \rg از یک درخواست در لحظه به $48$ درخواست هم‌زمان & \rg $150$ تا $220$ سند در دقیقه، صفر شکست در $1{,}996$ فراخوانی \\
\hline
\rg ازسرگیری & \rg فایل کناری برای سندهای بدون رابطه & \rg هیچ فراخوانی تکراری در ازسرگیری \\
\hline
\rg بستار و ارزیاب & \rg از حافظه به پردازش جریانی با اندیس بایتی & \rg فایل $1.5$ میلیون سطری بدون رشد حافظه \\
\hline
\rg بارگذار & \rg از چهار پرس‌وجو برای هر رابطه به دسته‌های $2{,}000$ تایی & \rg $12{,}437$ رابطه در $17.5$ ثانیه \\
\hline
\end{tabular}
\end{table}

درستی نسخه جریانی را با بازتولید دقیق عدد $28.64$ روی همان داده تأیید کردیم؛ اگر پردازش جریانی چیزی را جا می‌انداخت، این عدد جابه‌جا می‌شد.

\section{پشته فناوری}
\label{sec:stack}

زبان پیاده‌سازی \lr{Python} است. متن از \lr{PDF} و \lr{DOCX} با \lr{PyMuPDF} و \lr{python-docx} بیرون کشیده می‌شود و متن فارسی با \lr{hazm} نرمال می‌شود. در مرحله یک و دو از مدل \lr{gpt-4.1-mini} با دمای صفر و در مرحله سوم از مدل \lr{Claude Sonnet} استفاده شد. اعتبارسنجی خروجی با \lr{Pydantic} و تلاش دوباره با \lr{tenacity} انجام می‌شود، گراف در \lr{Neo4j 5} روی \lr{Docker} می‌نشیند و رابط وب با \lr{FastAPI} و \lr{vis-network} ساخته شده است.
